% Generated draft: requires mathematical and source review.
\documentclass[11pt]{article}
\usepackage[margin=25mm]{geometry}
\usepackage{fontspec}
\setmainfont{Latin Modern Roman}
\usepackage{amsmath,amssymb,mathtools,unicode-math}
\setmathfont{Latin Modern Math}
\usepackage{enumitem,needspace,xurl,hyperref}
\setlength{\parindent}{0pt}
\setlength{\parskip}{5pt}
\setlength{\emergencystretch}{3em}
\newcommand{\N}{\mathbb{N}}
\newcommand{\Z}{\mathbb{Z}}
\newcommand{\Q}{\mathbb{Q}}
\newcommand{\R}{\mathbb{R}}
\newcommand{\C}{\mathbb{C}}
\newcommand{\F}{\mathbb{F}}
\newcommand{\A}{\mathbb{A}}
\newcommand{\PP}{\mathbb{P}}
\newcommand{\E}{\mathbb{E}}
\newcommand{\eps}{\varepsilon}
\newcommand{\dd}{\,\mathrm{d}}
\begin{document}
\section*{UnsolvedMath ID: 20001286}
\textbf{Source problem code:} catalogue entry
\subsection{Definitions and mathematical statement}
S. Reisner's source problem asks to quantify the shadow movement theorem, citing [MR], without specifying a deficit, distance, normalization or remainder. In the Meyer--Reisner interpretation used by the attached aid, a shadow system moves generating points linearly in one fixed direction and takes their convex hull; the relevant convex functional is inverse Santal\'o-polar volume, not the volume product itself.

The aid supplies the planar benchmark
\[
K_t=\operatorname{conv}\{\pm(1+t,1),\ \pm(0,1),\ \pm(-1,0)\},\qquad-1<t<1.
\]
Only the first opposite pair moves, in the horizontal direction with scalar speeds $1,-1$; thus maximum generator speed is one. Central symmetry makes its Santal\'o point the origin, so its Santal\'o polar is the usual $K_t^\circ$. The exact areas and functional are
\[
|K_t|=3+t,\qquad |K_t^\circ|=3-t,\qquad F(t)=|K_t^\circ|^{-1}=\frac1{3-t}.
\]
For $s,u\in(-1,1)$, $0\le\lambda\le1$ and $m_\lambda=(1-\lambda)s+\lambda u$, define the Jensen deficit $J=(1-\lambda)F(s)+\lambda F(u)-F(m_\lambda)$. The aid states
\[
J=\frac{\lambda(1-\lambda)(u-s)^2}{(3-s)(3-u)(3-m_\lambda)}
\ge\frac{\lambda(1-\lambda)(u-s)^2}{64}.
\]
In particular the midpoint deficit is at least $(u-s)^2/256$, with zero at coincident parameters or endpoint weights. The contrasting volume product is $|K_t||K_t^\circ|=9-t^2$, concave in $t$, so it must not be substituted for the inverse-polar functional.

The aid also notes that parameter rescaling changes any parameter-based quadratic coefficient, translations may give genuine motion with zero deficit, and affine Meyer--Reisner equality families supply further null directions. A general positive quantitative remainder must therefore state speed/scale normalization and how affine null directions are removed. The planar formula is a test case, not a quantitative theorem for every shadow system.
\subsection{Short English statement}
Quantify shadow-system convexity with a specified normalization and treatment of affine equality directions. The supplied aid gives an exact unit-generator-speed hexagon deficit and warns that the volume product is a different, here concave, functional.
\subsection{Sources}
\begin{itemize}
\item Local UnsolvedMath registry, permanent ID $20001286$, source code catalogue entry; canonical \texttt{aim-convex-geometry-notes.json}, \texttt{notes[17]}.
\item American Institute of Mathematics, ``Mahler's conjecture and duality in convex geometry,'' stored problem 2 attributed to S. Reisner: \url{https://aimath.org/WWN/mahlerduality/mahlerduality.pdf}. The source cites [MR] without a bibliographic entry; a terminal page-number spill is not part of the question. No PDF inspection is claimed here.
\item Attached research attempt 1: partial result, recorded run status open, rigor proved, novelty confidence low; registry status partially solved, classification PARTIAL-PROGRESS. It identifies the Meyer--Reisner functional and supplies the model computation but no full theorem citation. This machine-generated, non-peer-reviewed aid requires independent verification before citation.
\end{itemize}
\end{document}
